% ═══════════════════════════════════════════════════════════════
% MT-01b-ML: Musterlösung Minitest Vorstellen + ser
% Spanisch Klasse 7 - Sequenz 1: Empezamos
% ═══════════════════════════════════════════════════════════════

\documentclass[11pt, a4paper]{article}

\usepackage[utf8]{inputenc}
\usepackage[T1]{fontenc}
\usepackage[ngerman]{babel}
\usepackage{helvet}
\renewcommand{\familydefault}{\sfdefault}

\usepackage[margin=2cm]{geometry}
\usepackage{enumitem}
\usepackage{tabularx}
\usepackage{booktabs}
\usepackage{xcolor}
\usepackage{fancyhdr}
\usepackage{lastpage}

% FARBEN
\definecolor{spanisch}{HTML}{c41e3a}
\definecolor{grau}{HTML}{888888}
\definecolor{hellgrau}{HTML}{f5f5f5}
\definecolor{loesung}{HTML}{c62828}

\pagestyle{fancy}
\fancyhf{}
\renewcommand{\headrulewidth}{0.4pt}
\renewcommand{\footrulewidth}{0.4pt}

\lhead{Spanisch -- Klasse 7}
\chead{\textbf{MT-01b -- MUSTERLÖSUNG}}
\rhead{\today}

\lfoot{10 Minuten}
\cfoot{}
\rfoot{Seite \thepage\ von \pageref{LastPage}}

\newcommand{\punktebox}[1]{\hfill\fbox{\small \_\_\_/#1}}
\newcommand{\luecke}[1]{\rule{#1}{0.4pt}}
\newcommand{\es}[1]{\textcolor{spanisch}{\textbf{#1}}}
\newcommand{\lsg}[1]{\textcolor{loesung}{\textbf{#1}}}

\begin{document}

% ═══════════════════════════════════════════════════════════════
% KOPFBEREICH
% ═══════════════════════════════════════════════════════════════

\begin{center}
    {\Large\bfseries\textcolor{spanisch}{Minitest: Vorstellen + ser}}

    \vspace{5pt}
    {\large\textcolor{loesung}{-- MUSTERLÖSUNG --}}
\end{center}

\vspace{10pt}

% ═══════════════════════════════════════════════════════════════
% AUFGABE 1: Konjugation ser (8 BE)
% ═══════════════════════════════════════════════════════════════

\noindent\textbf{\textcolor{spanisch}{Aufgabe 1 -- Das Verb ser konjugieren}} \punktebox{8}

\vspace{10pt}

\begin{tabularx}{\textwidth}{|l|X|}
\hline
\textbf{Pronomen} & \textbf{ser} \\
\hline
yo & \lsg{soy} \hfill \textit{(2 BE)} \\
\hline
tú & \lsg{eres} \hfill \textit{(2 BE)} \\
\hline
él / ella & \lsg{es} \hfill \textit{(2 BE)} \\
\hline
nosotros & \lsg{somos} \hfill \textit{(2 BE)} \\
\hline
\end{tabularx}

\vspace{5pt}
\textit{Hinweis: Die Formen von ser sind unregelmäßig und müssen auswendig gelernt werden.}

\vspace{15pt}

% ═══════════════════════════════════════════════════════════════
% AUFGABE 2: Lückentext ser (8 BE)
% ═══════════════════════════════════════════════════════════════

\noindent\textbf{\textcolor{spanisch}{Aufgabe 2 -- ser einsetzen}} \punktebox{8}

\vspace{10pt}

\begin{enumerate}[label=\alph*)]
    \item Yo \lsg{soy} de Alemania. \hfill \textit{(2 BE)}
    \item Tú \lsg{eres} estudiante. \hfill \textit{(2 BE)}
    \item Ella \lsg{es} de España. \hfill \textit{(2 BE)}
    \item Nosotros \lsg{somos} de Rostock. \hfill \textit{(2 BE)}
\end{enumerate}

\vspace{5pt}
\textit{Typischer Fehler: Verwechslung von ``soy'' (1. Person) und ``es'' (3. Person).}

\vspace{15pt}

% ═══════════════════════════════════════════════════════════════
% AUFGABE 3: Redemittel (6 BE)
% ═══════════════════════════════════════════════════════════════

\noindent\textbf{\textcolor{spanisch}{Aufgabe 3 -- Redemittel zuordnen}} \punktebox{6}

\vspace{10pt}

\begin{enumerate}[label=\arabic*.]
    \item \es{Me llamo...} = \lsg{a) Ich heiße...}
    \item \es{Soy de...} = \lsg{b) Ich bin aus...}
    \item \es{¿Cómo te llamas?} = \lsg{c) Wie heißt du?}
    \item \es{¿De dónde eres?} = \lsg{d) Woher bist du?}
\end{enumerate}

\vspace{5pt}
\textit{Bewertung: 4 richtig = 6 BE, 3 richtig = 5 BE, 2 richtig = 3 BE, 1 richtig = 2 BE}

\vspace{15pt}

% ═══════════════════════════════════════════════════════════════
% AUFGABE 4: Dialog (6 BE)
% ═══════════════════════════════════════════════════════════════

\noindent\textbf{\textcolor{spanisch}{Aufgabe 4 -- Dialog vervollständigen}} \punktebox{6}

\vspace{10pt}

\noindent
\textbf{A:} ¡Hola! ¿Cómo te llamas?

\vspace{0.4em}
\noindent
\textbf{B:} \lsg{Me llamo [Name]. / Soy [Name].} \hfill \textit{(2 BE)}

\vspace{0.6em}
\noindent
\textbf{A:} ¿De dónde eres?

\vspace{0.4em}
\noindent
\textbf{B:} \lsg{Soy de [Ort/Land]. / De [Ort/Land].} \hfill \textit{(2 BE)}

\vspace{0.6em}
\noindent
\textbf{A:} ¡Mucho gusto!

\vspace{0.4em}
\noindent
\textbf{B:} \lsg{¡Igualmente! / ¡Mucho gusto! / ¡Encantado/a!} \hfill \textit{(2 BE)}

\vspace{5pt}
\textit{Akzeptierte Antworten: Jede sinnvolle Antwort mit korrekter Struktur. Bei Antwort 3 auch: ``Yo también'' oder einfach ``Hola''.}

% ═══════════════════════════════════════════════════════════════
% PUNKTEÜBERSICHT
% ═══════════════════════════════════════════════════════════════

\vspace{20pt}
\hrule
\vspace{10pt}

\noindent
\begin{tabularx}{\textwidth}{|X|c|c|c|c||c|}
\hline
\textbf{Aufgabe} & 1 & 2 & 3 & 4 & \textbf{Gesamt} \\
\hline
\textbf{max. Punkte} & 8 & 8 & 6 & 6 & \textbf{28} \\
\hline
\end{tabularx}

\vspace{15pt}

\noindent\textbf{Notenskala (14-NP):}

\vspace{5pt}

\begin{tabular}{|c|c|c|c|c|c|c|c|c|c|c|c|c|c|c|}
\hline
NP & 14 & 13 & 12 & 11 & 10 & 9 & 8 & 7 & 6 & 5 & 4 & 3 & 2 & 1 \\
\hline
\% & 100 & 96 & 91 & 86 & 80 & 73 & 67 & 60 & 53 & 47 & 40 & 33 & 27 & 20 \\
\hline
BE & 28 & 27 & 25 & 24 & 22 & 20 & 19 & 17 & 15 & 13 & 11 & 9 & 8 & 6 \\
\hline
\end{tabular}

\end{document}
